% Wiederverwendbare Darstellung fuer Gegnerprofile

\definecolor{hpEnemy}{HTML}{F2EEE6}
\definecolor{hpEnemyLine}{HTML}{8A6A49}
\definecolor{hpEnemyInner}{HTML}{FBF9F4}

% Ein Gegnerprofil bildet einen geschlossenen Almanach-Eintrag.
\newtcolorbox{enemyprofile}[2][]{
  hpbox,
  title={#2},
  colback=hpEnemy,
  colframe=hpEnemyLine,
  boxed title style={colback=hpEnemyLine!22},
  fonttitle=\HPHeadingFont\large,
  #1
}

% Kompakte Teilbereiche innerhalb eines Gegnerprofils.
\newtcolorbox{enemysectionbox}[1]{
  enhanced,
  colback=hpEnemyInner,
  colframe=hpLine,
  boxrule=0.55pt,
  arc=1.5mm,
  left=3mm,
  right=3mm,
  top=2.5mm,
  bottom=2.5mm,
  before skip=0pt,
  after skip=0pt,
  title={#1},
  fonttitle=\HPUIFont\bfseries,
  fontupper=\HPUIFont,
  coltitle=hpBrown,
  attach title to upper={\quad}
}

% Werteband direkt unter dem Gegnernamen.
\newenvironment{enemystats}
  {\begingroup\HPUIFont\renewcommand{\arraystretch}{1.15}\setlength{\tabcolsep}{4pt}%
   \tabularx{\linewidth}{@{}>{\bfseries}l X >{\bfseries}l X >{\bfseries}l X@{}}}
  {\endtabularx\endgroup}

% Zweispaltiger Inhaltsbereich fuer Angriffe, KI und Sonderregeln.
\newenvironment{enemygrid}
  {\begin{tcbraster}[
      raster columns=2,
      raster column skip=4mm,
      raster row skip=3mm,
      raster valign=top
    ]}
  {\end{tcbraster}}

\newcommand{\enemyrole}[1]{\textbf{Rolle:} #1}
\newcommand{\enemydesign}[1]{\textbf{Designziel:} #1}
